% Figure (Chapter 3 v2): the methodology as a compact flowchart. One row
% per stage (common setup, characterization, development in two rows,
% evaluation); every step is described in Table~\ref{tab:steps}.
\begin{figure}[!htbp]
\centering
\begin{adjustbox}{max width=\linewidth}
\begin{tikzpicture}[
  font=\figfont,
  wf/.style={execute at begin node={\hyphenpenalty=10000\relax}, draw, rounded corners=2pt, fill=white, text width=2.55cm, align=center, inner sep=3pt, minimum height=1.3cm},
  wfdec/.style={wf, draw=newLine, very thick, fill=newFill},
  arr/.style={-{Stealth[length=1.8mm]}, thick},
  lane/.style={font=\figfont\bfseries, align=left, text width=2.0cm, anchor=west}
]
\def\wfxa{3.3}\def\wfxb{6.4}\def\wfxc{9.5}\def\wfxd{12.6}\def\wfxe{15.7}
\node[wf] (s1)  at (\wfxa,0)     {\textbf{1. Prepare data}\\ splits 100\,/\,200\,/\,200; LibriSpeech-PC};
\node[wf] (s2)  at (\wfxb,0)     {\textbf{2. Prompt conditions}\\ clean and seven degradations};
\node[wf] (s3)  at (\wfxc,0)     {\textbf{3. Check pipeline}\\ seeds, identities, pass counts};

\node[wf] (s4)  at (\wfxa,-1.9)  {\textbf{4. Run systems}\\ references and both payloads};
\node[wf] (s5)  at (\wfxb,-1.9)  {\textbf{5. Reproduction check}\\ against published numbers};
\node[wf] (s6)  at (\wfxc,-1.9)  {\textbf{6. Fix floors}\\ $s_{\min}$, $u_{\min}$ from the baseline};

\node[wf] (s7)  at (\wfxa,-3.8)  {\textbf{7. Formulate method}\\ $Q$, $r$, policy, proofs};

\node[wf] (s8)  at (\wfxa,-5.7)  {\textbf{8. Fix protocol}\\ $\alpha$, $\delta$, grid of $\lambda_r$};
\node[wf] (s9)  at (\wfxb,-5.7)  {\textbf{9. Calibrate}\\ binomial tests in order};
\node[wf] (s10) at (\wfxc,-5.7) {\textbf{10. Threshold certified?}\\ no: use DMOSpeech~2};

\node[wf] (s11) at (\wfxa,-7.6)  {\textbf{11. Check the bound}\\ test risk, 100 re-splits};
\node[wf] (s12) at (\wfxb,-7.6)  {\textbf{12. Run compared systems}\\ test split, LibriSpeech-PC};
\node[wf] (s13) at (\wfxc,-7.6)  {\textbf{13. Compare}\\ cost, quality, variety};
\node[wf] (s14) at (\wfxd,-7.6)  {\textbf{14. Robustness}\\ conditions, $n$, $\alpha$};
\node[wf, fill=frozenFill] (s15) at (\wfxe,-7.6) {\textbf{15. Report}\\ including negative results};

\foreach \a/\b in {s1/s2,s2/s3,s4/s5,s5/s6,s8/s9,s9/s10,s11/s12,s12/s13,s13/s14,s14/s15}
  \draw[arr] (\a) -- (\b);
\draw[arr] (s3.south)  -- (\wfxc,-0.95) -- (\wfxa,-0.95) -- (s4.north);
\draw[arr] (s6.south)  -- (\wfxc,-2.85) -- (\wfxa,-2.85) -- (s7.north);
\draw[arr] (s7.south)  -- (s8.north);
\draw[arr] (s10.south) -- (\wfxc,-6.65) -- (\wfxa,-6.65) -- (s11.north);

\begin{scope}[on background layer]
  \foreach \y in {0,-1.9,-3.8,-5.7,-7.6}
    \draw[draw=frozenLine, fill=frozenFill, rounded corners=4pt] (-0.45,\y+0.75) rectangle (17.2,\y-0.75);
\end{scope}
\node[lane] at (-0.3,0)    {Common setup\\[1pt]\mdseries all splits};
\node[lane] at (-0.3,-1.9) {Objective 1\\[1pt]\mdseries development};
\node[lane] at (-0.3,-3.8) {Objective 2\\[1pt]\mdseries formulate\\no data};
\node[lane] at (-0.3,-5.7) {Objective 2\\[1pt]\mdseries calibrate\\calibration};
\node[lane] at (-0.3,-7.6) {Objective 3\\[1pt]\mdseries test};
\end{tikzpicture}
\end{adjustbox}
\caption{Methodology of the study. Each shaded row is one stage, labelled with the data it uses; Table~\ref{tab:steps} describes every step. All networks stay frozen throughout.}
\label{fig:workflow}
\end{figure}
